\documentclass{article}
\usepackage[utf8]{inputenc}
\title{Premium AI Mastery Course: Prompt Engineering Fundamentals}
\author{Mandarinii}
\date{2026-03-17}
\begin{document}
\maketitle
\section{Introduction}
In this lesson, we explore the fundamentals of prompt engineering, which is essential for mastering AI interactions.
\section{The 5-Block Formula}
The 5-block formula is an effective method for crafting prompts. Each block focuses on a specific aspect:
\begin{enumerate}
    \item \textbf{Role}: Define the role of the AI.
    \item \textbf{Task}: Specify what the AI needs to accomplish.
    \item \textbf{Context}: Provide background information necessary for the task.
    \item \textbf{Constraints}: Set any limitations or guidelines.
    \item \textbf{Format}: Indicate the desired format of the output.
\end{enumerate}
\section{Examples}
\subsection{Example 1: Chatbot}
\textbf{Role:} You are a helpful chatbot.\newline
\textbf{Task:} Answer questions about gardening.\newline
\textbf{Context:} User is new to gardening and needs beginner tips.\newline
\textbf{Constraints:} Responses should be concise and easy to understand.\newline
\textbf{Format:} Provide a bullet-point list of tips.
\subsection{Example 2: Content Writer}
\textbf{Role:} You are a creative content writer.\newline
\textbf{Task:} Write a blog post about AI advancements.\newline
\textbf{Context:} The audience is tech enthusiasts.\newline
\textbf{Constraints:} Limit to 600 words and include at least three references.\newline
\textbf{Format:} Standard blog post format.
\section{Practical Templates}
\subsection{General Prompt Template}
\begin{verbatim}
You are a {Role}. Your task is to {Task}. The context is as follows: {Context}. Please follow these constraints: {Constraints}. Format the output as {Format}.
\end{verbatim}
\section{Conclusion}
Following the 5-block formula enables more effective communication with AI, enhancing the overall experience. As you practice, tailor prompts to fit specific needs and contexts.
\end{document}